\newpage
\reviewsection{Reflection on the Lesson}

The strongest part of the lesson was its connection to ordinary classroom behavior. Students often understand that they should be ``nice,'' but that word is too general to guide them during noise, frustration, partner work, login confusion, unequal pacing, or group cleanup. The lesson gave those moments usable language: help without taking over, wait without rushing, ask without assuming, respect without isolating, and include without forcing.

The placement evidence made the reflection more honest. In technology-rich classrooms, competence can be hidden by login friction, unclear click paths, reading demands, sensory load, role ambiguity, or the pace of a partner. A student who pauses, repeats a question, moves away, watches before joining, or refuses a material may be communicating confusion, overload, uncertainty, or a need for more predictable access. Those reactions should not automatically be interpreted as inability, disrespect, or lack of interest.

\begin{figure}[htbp]
\centering
\resizebox{0.9\textwidth}{!}{\begin{tikzpicture}[center/.style={ellipse,draw=PrimaryRed,very thick,fill=SoftGray,minimum width=2.6in,minimum height=1in,align=center},need/.style={rounded corners=3mm,draw=black!55,thick,minimum width=2.15in,minimum height=.78in,align=center},goal/.style={rounded corners=3mm,draw=DeepPurple,very thick,fill=SoftLavender,minimum width=3.2in,minimum height=.9in,align=center},arrow/.style={-{Latex[length=3mm]},thick,draw=black!65}]
\node[center] (c) {\textbf{Participation is shaped by access}};
\node[need,fill=SoftBlue,above left=10mm and 16mm of c] (a) {\textbf{Communication}\\speech, AAC, gesture, writing, silence};
\node[need,fill=SoftGreen,above right=10mm and 16mm of c] (b) {\textbf{Sensory access}\\noise, space, movement, breaks};
\node[need,fill=SoftGold,below left=10mm and 16mm of c] (d) {\textbf{Processing and pacing}\\wait time, sequence, predictability};
\node[need,fill=SoftLavender,below right=10mm and 16mm of c] (e) {\textbf{Peer belonging}\\invitation, consent, roles, respect};
\node[goal,below=27mm of c] (g) {Access $\rightarrow$ dignity $\rightarrow$ agency $\rightarrow$ belonging};
\draw[arrow] (a)--(c); \draw[arrow] (b)--(c); \draw[arrow] (d)--(c); \draw[arrow] (e)--(c); \draw[arrow] (c)--(g);
\end{tikzpicture}}
\caption{Neurodivergent access map. Participation is interpreted through communication, sensory access, processing and pacing, and peer belonging rather than through behavior alone.}
\label{fig:access-map}
\end{figure}

Figure~\ref{fig:access-map} captures the central reflection: behavior cannot be understood apart from access. A delayed answer may indicate processing time or another communication mode. Moving away may indicate sensory load. Difficulty joining a coding task may reflect unclear steps rather than low ability. A quiet or new student outside the group may need an explicit invitation and a safe role rather than pressure.

The lesson also showed me that peer culture is part of positive behavioral support. If the classroom teaches students to invite, wait, ask, respect, and include, support does not rest only on the teacher or on the autistic student. Peers learn that access is something the whole classroom helps create.

The part I would revise is pacing. Ten minutes was enough for one core idea and one realistic scenario, but not enough for deeper discussion of communication differences, sensory needs, privacy, consent, and the range of autistic experiences. I would revisit one action before each partner-heavy activity so the language becomes part of classroom routines rather than a one-time presentation.

\reviewsection{Puzzle Plan and EQUITAS Alignment}

The \textbf{Puzzle Plan} is my interdisciplinary framework for combining lived experience, educational practice, classroom evidence, data, and research instead of relying on one isolated explanation. \textbf{EQUITAS} is the equity-focused lens within that framework; it asks whose knowledge is included, what context is missing, and where power or bias may shape a decision.

Through this lens, no diagnosis, behavior, observation, or peer comment should stand alone. The stronger question is what communication option, wait time, sensory support, relationship move, or environmental change might increase access. The goal is not pity, charity, or visible compliance. The goal is dignity, agency, participation, and belonging.
